% ===================================================================
%  图表第 13 号 — 旅馆。— 饭店。— 咖啡馆。
%  Traduction 2026 d'apres texto/io, controlee sur texto/fr, texto/en
%  et texto/es. Consigne : voir texto/zh/10-tubiao-01.tex.
%
%  Deux notes, marquees toutes deux « (*) », et deux appels : celui du
%  bloc 01-2 (la chambre) et celui du bloc 09-1 (le menu). L'ordre les
%  departage.
% ===================================================================

%%K t13-noto-1 noto
\VUnotes{143.2mm}{%
\textit{\VUgras{(*) 卧房的细目见图表第 6 号。}}%
}

%%K t13-apar-1 apar
\VUsaut{3.0mm}
\VUtitre{15.0pt}{60}{图表第 13 号}
\VUsaut{1.6mm}
\VUpk{10.2pt}{40}{旅馆。--- 饭店。}
\VUsaut{2.0mm}
\VUpk{10.2pt}{40}{咖啡馆。}
\VUsaut{3.4mm}

%%K t13-01-1 p
1. --- 昨天傍晚到了鲁瓦扬，我住进\textit{海洋旅馆}，那是一个朋友介绍的。

%%K t13-01-2 p
他们给了我一间很好的\VUgras{房间} \textsuperscript{(2)}，在三
\VUgras{层} \textsuperscript{(1)} 上，我在那里睡得很好(*)。

%%K t13-02-1 p
2. --- 今天早上出了房间——\VUgras{女仆} \textsuperscript{(3)} 已经把早点送到
那里——我走进\VUgras{吸烟室} \textsuperscript{(5)}，那里同时也是\VUgras{阅览室}
\textsuperscript{(4)}，一面翻\VUgras{报} \textsuperscript{(6)}，一面抽\VUgras{纸烟}
\textsuperscript{(8)}。看完了报，我坐\VUgras{电梯} \textsuperscript{(9)} 下楼，出去在
城里散步。

%%K t13-03-1 p
3. --- 因为忘了问旅馆的\VUgras{账房} \textsuperscript{(12)} 几点开
\VUgras{公共餐桌} \textsuperscript{(11)}，我早早回来，趁便在旅馆的\VUgras{浴室} \textsuperscript{(10)}
里洗了个澡。然后我到\VUgras{账台} \textsuperscript{(15)} 去给一位朋友打电话。
可是\VUgras{电话} \textsuperscript{(13)} 占着线；\VUgras{女账房} \textsuperscript{(14)}
看我为难——她正把一张\VUgras{账单} \textsuperscript{(17)} 记进\VUgras{旅客簿}
\textsuperscript{(16)}——就请\VUgras{门房} \textsuperscript{(18)} 打发\VUgras{侍应生}
\textsuperscript{(19)} 骑他的\VUgras{自行车} \textsuperscript{(20)} 去替我跑一趟。

%%K t13-04-1 p
4. --- 我谢了她，出去给\VUgras{脚夫} \textsuperscript{(24)}（就是那个打杂的）
一点赏钱；他用他的\VUgras{手推车} \textsuperscript{(25)} 从车站给我运来了
\VUgras{帽盒} \textsuperscript{(26)}、\VUgras{皮箱} \textsuperscript{(27)} 和
\VUgras{旅行袋} \textsuperscript{(28)}。\VUgras{邮差} \textsuperscript{(21)} 正好这时
到了，交给我一封\VUgras{信} \textsuperscript{(22)}。

%%K t13-05-1 p
5. --- 我在门口\VUgras{信箱} \textsuperscript{(23)} 旁边又待了一会儿，看街上的
来往。\VUgras{车夫} \textsuperscript{(30)} 是\VUgras{公共马车}
\textsuperscript{(29)} 的，他正把一只\VUgras{大箱} \textsuperscript{(31)} 装上车，
那是一位\VUgras{客人} \textsuperscript{(32)} 的，\VUgras{旅馆老板} \textsuperscript{(33)} 在同那客人道别。一个年轻的
\VUgras{电报员} \textsuperscript{(35)} 一点也不急，慢慢送来一封给旅馆某位客人的
\VUgras{电报} \textsuperscript{(34)}；一位\VUgras{游客} \textsuperscript{(36)} 肩上挎着
\VUgras{照相机} \textsuperscript{(38)}，正在查他的\VUgras{指南} \textsuperscript{(37)}。

%%K t13-tit-1 sub
\VUsaut{4.0mm}
\VUpk{10.2pt}{40}{午饭时候。}
\VUsaut{3.0mm}

%%K t13-06-1 p
6. --- 栅栏前面，一个安闲的\VUgras{跑腿的} \textsuperscript{(39)} 在他的
\VUgras{挑钩} \textsuperscript{(40)} 和\VUgras{擦鞋箱} \textsuperscript{(41)} 中间打盹；
一个年轻的\VUgras{洗衣妇} \textsuperscript{(42)} 提着一\VUgras{篮} \textsuperscript{(43)}
衣裳，笑着看门口那位\VUgras{领班} \textsuperscript{(44)} 一本正经的样子。

%%K t13-07-1 p
7. --- \VUgras{厨子} \textsuperscript{(45)} 从\VUgras{厨房} \textsuperscript{(47)}
里出来——厨房在\VUgras{地下室} \textsuperscript{(48)}——去打发一个
\VUgras{小厨工} \textsuperscript{(46)} 办事；我一看见就想起午饭时候近了，便到院子那边的饭店去。那里一个
\VUgras{开汽车的} \textsuperscript{(50)} 正停他的\VUgras{汽车} \textsuperscript{(49)}，
一个\VUgras{机匠} \textsuperscript{(52)} 照着\VUgras{骑车人} \textsuperscript{(53)} 的
指点，在修一辆刚出了事的\VUgras{汽油三轮} \textsuperscript{(51)}。

%%K t13-08-1 p
8. --- 我问一个年轻的\VUgras{侍役} \textsuperscript{(54)}，他端着
\VUgras{几杯啤酒} \textsuperscript{(55)}，公共餐桌是不是摆在凉篷底下；他说是，我就在一张桌子旁
坐下，吃了一顿很好的饭。

%%K t13-noto-2 noto
\VUnotes{143.2mm}{%
\textit{\VUgras{(*) 借着词汇表里那份菜名单子，老师很容易把下面的菜单变换
花样。}}%
}

%%K t13-tit-2 sub
\VUsaut{4.0mm}
\VUpk{10.2pt}{40}{一顿很好的午饭。}
\VUsaut{3.0mm}

%%K t13-09-1 p
9. --- 侍者送来午饭的菜单(*)和酒单。我点了菜，\VUgras{冷盘}要了
\VUgras{黄油}\VUgras{对虾}，头一道菜要了\VUgras{卡昂式牛肚}。\VUgras{鱼}
我没要，只要了一份羊\VUgras{腿}，\VUgras{蔬菜}要了奶油\VUgras{菠菜}。
后来那些\VUgras{甜食}我一样也不中意，就叫他们送上各样点心：\VUgras{奶酪}、
\VUgras{水果}和\VUgras{饼干}。我又添了一瓶很好的圣埃米利翁\VUgras{酒}；
这顿饭吃得很满意，我给了\VUgras{侍者} \textsuperscript{(57)} 一笔可观的
\VUgras{小费} \textsuperscript{(59)}，就离了桌。

%%K t13-10-1 p
10. --- \VUgras{经理} \textsuperscript{(60)} 看我要走，请我到阳台上去看看
\VUgras{大海} \textsuperscript{(62)} 那壮丽的全景。远处看得出小小的打鱼
\VUgras{船} \textsuperscript{(63)}、\VUgras{帆船} \textsuperscript{(64)} 和一艘巨大的
\VUgras{邮轮} \textsuperscript{(65)}，它黑黑的影子映在白\VUgras{云} \textsuperscript{(67)}
上，云在\VUgras{天边} \textsuperscript{(66)}。一阵轻风吹动着那面\VUgras{旗}
\textsuperscript{(70)}，旗插在\VUgras{招牌} \textsuperscript{(69)} 上面，那是
\VUgras{门厅} \textsuperscript{(68)} 的招牌。

%%K t13-tit-3 sub
\VUsaut{3.0mm}
\VUpk{10.2pt}{40}{消遣。}
\VUsaut{3.0mm}

%%K t13-11-1 p
11. --- 这景致太有意思了，我打定主意第二天上咖啡馆的露台去，带一副
\VUgras{观剧镜} \textsuperscript{(71)} 或者一架\VUgras{望远镜} \textsuperscript{(72)}。

%%K t13-12-1 p
12. --- 离了阳台，我到旅馆的咖啡馆去喝一杯\VUgras{饮料} \textsuperscript{(73)}，
打一局\VUgras{台球} \textsuperscript{(75)}。我一直穿过咖啡厅，那里许多客人在下
\VUgras{象棋} \textsuperscript{(78)}、下\VUgras{跳棋} \textsuperscript{(79)}、打
\VUgras{多米诺} \textsuperscript{(80)}、玩\VUgras{十五子} \textsuperscript{(81)} 或者打
\VUgras{纸牌} \textsuperscript{(82)}；我径直上了\VUgras{台球室} \textsuperscript{(74)}。

%%K t13-13-1 p
13. --- 打完球，我下到底层，向侍者要了一个\VUgras{信封} \textsuperscript{(84)}、
一些\VUgras{信纸} \textsuperscript{(83)}、一张\VUgras{邮票} \textsuperscript{(85)}、
一支\VUgras{笔} \textsuperscript{(87)} 和一点\VUgras{墨水} \textsuperscript{(86)}，好写
一封信。

%%K t13-13-2 p
然后我叫了一个\VUgras{车夫} \textsuperscript{(92)}，他的\VUgras{马车}
\textsuperscript{(88)} 就停在咖啡馆门前。我问他按钟点和按趟怎么算；讲好了，他
给我打开\VUgras{车门} \textsuperscript{(89)}，扶我坐进去。他又上了他的
\VUgras{驭座} \textsuperscript{(91)}，用\VUgras{鞭子} \textsuperscript{(93)} 麻利地
一抖，马就跑起来，我们便去兜了一趟极舒服的风。
\VUsaut{6.0mm}
\VUfilet{22mm}
